% ===== PRÉAMBULE CANONIQUE — à coller VERBATIM en tête de CHAQUE .tex =====
\documentclass[11pt,a4paper]{article}
\usepackage[utf8]{inputenc}\usepackage[T1]{fontenc}\usepackage{lmodern}
\usepackage[french]{babel}
\usepackage{amsmath,amssymb,mathtools}\usepackage{icomma}
\usepackage{siunitx}\sisetup{locale=FR}  % \qty{}{}, \si{}, \ang{}
\usepackage{xcolor}\usepackage{colortbl}
\usepackage{tikz}\usepackage{pgfplots}\pgfplotsset{compat=1.18}
\usepackage[siunitx,europeanresistors]{circuitikz} % circuits (élec.)
\usepackage[most]{tcolorbox}\usepackage{enumitem}\usepackage{array,booktabs}
\usepackage[a4paper,margin=2.2cm]{geometry}
\usetikzlibrary{arrows.meta,calc,angles,quotes,positioning,patterns,%
  backgrounds,shapes.geometric,decorations.pathmorphing,decorations.markings,intersections}
\definecolor{primary}{RGB}{222,49,99}\definecolor{primarydark}{RGB}{150,22,55}
\definecolor{defcol}{RGB}{0,114,178}\definecolor{methcol}{RGB}{52,92,148}
\definecolor{excol}{RGB}{0,135,90}\definecolor{attncol}{RGB}{200,40,55}
\tcbset{boxbase/.style={enhanced,breakable,boxrule=0.9pt,arc=2.5pt,
  left=3mm,right=3mm,top=2.3mm,bottom=2.3mm,fonttitle=\bfseries,coltitle=white}}
% --- boîtes TUTORIEL (noms EXACTS, ne pas renommer) ---
\newtcolorbox{cle}[1][]{boxbase,colback=defcol!6,colframe=defcol,
  title={\textsc{À retenir}\ifx\relax#1\relax\relax\else\textmd{ : \itshape #1}\fi}}
\newtcolorbox{methode}[1][]{boxbase,colback=methcol!7,colframe=methcol,sharp corners,
  title={\textsc{Méthode}\ifx\relax#1\relax\relax\else\textmd{ --- \itshape #1}\fi}}
\newtcolorbox{exemplebox}[1][]{boxbase,colback=excol!5,colframe=excol,
  title={\textsc{Exemple}\ifx\relax#1\relax\relax\else\textmd{ : \itshape #1}\fi}}
\newtcolorbox{piege}[1][]{boxbase,colback=attncol!6,colframe=attncol,
  title={\textsc{Piège}\ifx\relax#1\relax\relax\else\textmd{ : \itshape #1}\fi}}
\newtcolorbox{remarque}[1][]{boxbase,colback=black!4,colframe=black!45,
  title={\textsc{Remarque}\ifx\relax#1\relax\relax\else\textmd{ : \itshape #1}\fi}}
% --- boîtes EXERCICE / CORRIGÉ (noms EXACTS, auto-comptées) ---
\newtcolorbox[auto counter]{exo}[1][]{boxbase,colback=primary!4,colframe=primary,
  title={\textsc{Exercice}~\thetcbcounter\ifx\relax#1\relax\relax\else\textmd{ --- \itshape #1}\fi}}
\newtcolorbox[auto counter]{corrige}[1][]{boxbase,colback=excol!4,colframe=excol,
  title={\textsc{Corrigé}~\thetcbcounter\ifx\relax#1\relax\relax\else\textmd{ --- \itshape #1}\fi}}
\newcommand{\vv}[1]{\vec{#1}}\newcommand{\ds}{\displaystyle}
\newcommand{\R}{\mathbb{R}}\newcommand{\N}{\mathbb{N}}\newcommand{\Z}{\mathbb{Z}}
% =========================================================================
\begin{document}
\sisetup{per-mode=symbol}

\begin{exo}[PFD sur un plan incliné avec frottement]
Un bloc de masse $m=\qty{5}{\kilogram}$ glisse vers le bas d'un plan incliné à $\alpha=\ang{30}$.
Le frottement vaut $F_f=\qty{10}{\newton}$ ; $g=\qty{10}{\metre\per\second\squared}$, $\sin\ang{30}=\num{0,5}$.
\begin{enumerate}[label=\alph*)]
  \item Écris le PFD selon l'axe $Ox$ (le long du plan, sens de la descente).
  \item Calcule l'accélération du bloc.
\end{enumerate}
\end{exo}

\begin{exo}[le traîneau tiré horizontalement]
On tire un traîneau de masse $m=\qty{40}{\kilogram}$ sur une piste horizontale avec une force
horizontale $\vv F=\qty{200}{\newton}$. Le coefficient de frottement cinétique vaut $\mu_c=\num{0,3}$
($g=\qty{10}{\metre\per\second\squared}$).
\begin{enumerate}[label=\alph*)]
  \item Calcule la force de frottement $F_f=\mu_c\, mg$.
  \item En déduis l'accélération du traîneau.
\end{enumerate}
\end{exo}

\begin{exo}[la caisse à vitesse constante]
Une caisse de masse $m=\qty{100}{\kilogram}$, posée sur un sol horizontal, est tirée par une force
horizontale de \qty{100}{\newton} ; elle se déplace en ligne droite à \emph{vitesse constante}. On
indique que le coefficient de frottement cinétique vaut $\mu_c=\num{0,3}$ ($g=\qty{10}{\metre\per\second\squared}$).
\begin{enumerate}[label=\alph*)]
  \item Que vaut la force de frottement qui s'exerce sur la caisse ?
  \item La donnée $\mu_c=\num{0,3}$ était-elle nécessaire ? Explique.
\end{enumerate}
\end{exo}

\begin{exo}[coefficient de frottement par la distance d'arrêt]
Un palet est lancé à \qty{4}{\metre\per\second} sur une table horizontale. À cause du frottement, il
s'arrête après \qty{10}{\metre} ($g=\qty{10}{\metre\per\second\squared}$).
\begin{enumerate}[label=\alph*)]
  \item Calcule la décélération à partir de $v^2=2a\,d$.
  \item Sachant que $F_f=\mu_c\, mg=ma$, déduis-en le coefficient de frottement cinétique $\mu_c$.
\end{enumerate}
\end{exo}

\begin{exo}[monter une caisse le long d'une pente]
Une personne pousse une caisse de masse $m=\qty{10}{\kilogram}$ à \emph{vitesse constante} le long d'un
plan incliné, en exerçant une force $\vv F=\qty{200}{\newton}$ \emph{parallèle} à la pente. Le plan a une
hauteur de \qty{3}{\metre} pour une longueur de \qty{5}{\metre} (donc $\sin\alpha=\tfrac{3}{5}$) ;
$g=\qty{10}{\metre\per\second\squared}$.
\begin{enumerate}[label=\alph*)]
  \item La caisse monte à vitesse constante : quelle relation lie $F$, $mg\sin\alpha$ et le frottement $F_f$ ?
  \item Calcule l'intensité du frottement $F_f$.
\end{enumerate}
\end{exo}
\end{document}
